\documentclass[12pt, a4paper, oneside]{ctexart}
\usepackage{amsmath, amsthm, amssymb, bm, xcolor, tikz, framed, graphicx, hyperref, mathrsfs}
\usepackage{bookmark}
\usepackage{pgfplots}
\pgfplotsset{compat=1.18}

\newcommand{\bra}[1]{\langle #1 |}
\newcommand{\ket}[1]{| #1 \rangle}
\newcommand{\braket}[2]{\langle #1 | #2 \rangle}

\title{Quantum Mechanics Homework (Week 10)}
\author{Tang Yi \\ [学号已脱敏]}
\usepackage[a4paper, left=2cm, right=2cm, top=2cm, bottom=2cm]{geometry}
\linespread{1.3} % 行距更紧凑
\definecolor{shadecolor}{RGB}{241, 241, 255}

\newcounter{problemname}
\newenvironment{problem}
  {\begin{shaded}\stepcounter{problemname}\par\noindent\textbf{题目 \theproblemname. }}
  {\end{shaded}\par}
\newenvironment{solution}
  {\par\noindent\textbf{题目 \theproblemname 的解答. }}
  {\par}
\newenvironment{note}
  {\par\noindent\textbf{题目 \theproblemname 的注记. }}
  {\par}

\begin{document}

\maketitle

\begin{problem}
  考虑$\delta$势阱$V(x) =−\gamma \delta(x)$中质量为$m$的粒子，处于束缚本征态。分别用坐标表象积分和费曼海尔曼定理求动能和势能的期望值。
\end{problem}

\begin{solution}
  在势阱
$$
V(x)=-\gamma\delta(x)
$$
中，束缚态波函数可写为
$$
\psi(x)=A e^{-\kappa |x|}.
$$

由含 $\delta$ 势的导数跳跃条件可得
$$
\kappa=\frac{m\gamma}{\hbar^2},
\qquad
A=\sqrt{\kappa},
\qquad
E=-\frac{\hbar^2\kappa^2}{2m}
=-\frac{m\gamma^2}{2\hbar^2}.
$$

为求期望值：

1. 坐标表象积分法

势能期望值
$$
\langle V\rangle
=-\gamma|\psi(0)|^2
=-\gamma\kappa
=-\frac{m\gamma^2}{\hbar^2}.
$$

动能期望值
$$
\langle T\rangle
=\frac{\hbar^2}{2m}\int_{-\infty}^\infty |\psi’(x)|^2 dx
=\frac{\hbar^2}{2m}\kappa^2
=\frac{m\gamma^2}{2\hbar^2}.
$$

2. 费曼–海尔曼定理

对参数 $\gamma$，有
$$
\frac{\partial H}{\partial\gamma}=-\delta(x),
\qquad
E(\gamma)=-\frac{m\gamma^2}{2\hbar^2}.
$$

因此
$$
\frac{\mathrm{d}E}{\mathrm{d}\gamma}
=-\frac{m\gamma}{\hbar^2}
=-\langle\delta(x)\rangle,
$$
从而
$$
\langle\delta(x)\rangle
=\frac{m\gamma}{\hbar^2}.
$$

势能期望值为
$$
\langle V\rangle
=-\gamma\langle\delta(x)\rangle
=-\frac{m\gamma^2}{\hbar^2}.
$$

最终
$$
\langle T\rangle=\frac{m\gamma^2}{2\hbar^2},
\qquad
\langle V\rangle=-\frac{m\gamma^2}{\hbar^2},
\qquad
E=-\frac{m\gamma^2}{2\hbar^2}.
$$
\end{solution}

\begin{problem}
  一维谐振子处于状态：

  \centerline{$\psi(x) = \dfrac{1}{\sqrt{2}} \psi_0(x) + \dfrac{1}{\sqrt{2}} \psi_1(x)$}

  在海森堡绘景下求在$t$时刻的动量期望值。
\end{problem}

\begin{solution}
  设一维谐振子的能量本征态满足
\[
H\psi_n = E_n\psi_n,\qquad E_n=\hbar\omega\left(n+\frac12\right).
\]
给定初态
\[
\psi(x,0)=\frac{1}{\sqrt{2}}\psi_0(x)+\frac{1}{\sqrt{2}}\psi_1(x).
\]

在海森堡绘景中，算符随时间演化：
\[
\hat{p}(t)=e^{iHt/\hbar}\hat{p}\,e^{-iHt/\hbar}.
\]

谐振子中有关系
\[
\hat{x}=\sqrt{\frac{\hbar}{2m\omega}}(a+a^\dagger),\qquad
\hat{p}=i\sqrt{\frac{m\hbar\omega}{2}}(a^\dagger-a).
\]

利用
\[
a(t)=a\,e^{-i\omega t},\qquad a^\dagger(t)=a^\dagger e^{i\omega t},
\]
可得
\[
\hat{p}(t)=i\sqrt{\frac{m\hbar\omega}{2}}\left(a^\dagger e^{i\omega t}-a e^{-i\omega t}\right).
\]

动量期望值为
\[
\langle p(t)\rangle
=\langle\psi(0)|\hat{p}(t)|\psi(0)\rangle.
\]

代入初态：
\[
|\psi(0)\rangle=\frac1{\sqrt2}(|0\rangle+|1\rangle).
\]

只需用
\[
a|1\rangle=\sqrt{1}|0\rangle,\qquad 
a^\dagger|0\rangle=\sqrt{1}|1\rangle,
\]
其余零。

于是
\[
\langle p(t)\rangle
= i\sqrt{\frac{m\hbar\omega}{2}}
\frac12\left(
\langle1|a^\dagger|0\rangle e^{i\omega t}
-\langle0|a|1\rangle e^{-i\omega t}
\right).
\]

注意
\[
\langle1|a^\dagger|0\rangle=\langle0|a|1\rangle=1,
\]
因此
\[
\langle p(t)\rangle
= i\sqrt{\frac{m\hbar\omega}{2}}\,
\frac12\left(e^{i\omega t}-e^{-i\omega t}\right).
\]

化简：
\[
e^{i\omega t}-e^{-i\omega t}=2i\sin(\omega t),
\]
得到最终结果
\[
\boxed{
\langle p(t)\rangle
= \sqrt{\frac{m\hbar\omega}{2}}\,\sin(\omega t)
}.
\]

\end{solution}

\begin{problem}
  $D_x(a) = \exp \left[-a \frac{\partial}{\partial x}\right] = \exp \left[-i a \hat{p}_x / \hbar\right]$ 表示 $x$ 方向平移距离 $a$ 的算符。证明下列形式的波函数（称为布洛赫波函数）
$$
\psi(x) = e^{i k x} \phi_k(x), \quad \phi_k(x+a) = \phi_k(x)
$$
是 $D_x(a)$ 的本征态，相应本征值为 $e^{-i k a}$。
\end{problem}


\begin{solution}
  \[
D_x(a)=\exp\!\left[-a\frac{\partial}{\partial x}\right]
=\exp\!\left[-\frac{i}{\hbar}a\hat p_x\right].
\]
该算符对任意函数的作用可由指数算符的泰勒展开给出：
\[
D_x(a)f(x)=\sum_{n=0}^{\infty}\frac{(-a)^n}{n!}\frac{\partial^n f(x)}{\partial x^n}
=f(x-a),
\]
即 $D_x(a)$ 将函数平移距离 $a$。

考虑布洛赫形式的波函数
\[
\psi(x)=e^{ikx}\phi_k(x), \qquad \phi_k(x+a)=\phi_k(x),
\]
作用平移算符得到
\[
D_x(a)\psi(x)=\psi(x-a)
=e^{ik(x-a)}\phi_k(x-a).
\]
利用 $\phi_k$ 的周期性 $\phi_k(x-a)=\phi_k(x)$，有
\[
D_x(a)\psi(x)
=e^{ik(x-a)}\phi_k(x)
=e^{-ika}\,e^{ikx}\phi_k(x)
=e^{-ika}\psi(x).
\]
因此 $\psi(x)$ 为 $D_x(a)$ 的本征态，其本征值为$e^{-ika}$.
\end{solution}
\begin{problem}
  阱底范围在 $(0, a)$ 的无限深势阱中有两个质量为 $m$ 的粒子，设粒子之间没有相互作用，对 (1) 全同玻色子；(2) 全同费米子；(3) 可分辨粒子的情况，分别写出基态和第一激发态的本征能量、本征态和简并度。
\end{problem}
\begin{solution}
  设单粒子在 $(0,a)$ 无限深势阱中的本征函数与本征能量为
\[
\psi_n(x)=\sqrt{\frac{2}{a}}\sin\!\left(\frac{n\pi x}{a}\right),\qquad
E_n=\frac{n^2\pi^2\hbar^2}{2ma^2}.
\]

\noindent\textbf{（1）全同玻色子}

基态：
\[
\Psi_{\mathrm{gs}}(x_1,x_2)=\psi_1(x_1)\psi_1(x_2),\qquad
E_{\mathrm{gs}}=2E_1,
\]
简并度：$1$。

第一激发态：
\[
\Psi_{1}(x_1,x_2)=\frac{1}{\sqrt{2}}\big[\psi_1(x_1)\psi_2(x_2)+\psi_2(x_1)\psi_1(x_2)\big],
\qquad E_1=E_1+E_2,
\]
简并度：$1$。

\bigskip

\noindent\textbf{（2）全同费米子}

基态：
\[
\Psi_{\mathrm{gs}}(x_1,x_2)=\frac{1}{\sqrt{2}}\big[\psi_1(x_1)\psi_2(x_2)-\psi_2(x_1)\psi_1(x_2)\big],\qquad
E_{\mathrm{gs}}=E_1+E_2,
\]
简并度：$1$。

第一激发态：
\[
\Psi_1(x_1,x_2)=\frac{1}{\sqrt{2}}\big[\psi_1(x_1)\psi_3(x_2)-\psi_3(x_1)\psi_1(x_2)\big],\qquad
E_1=E_1+E_3,
\]
简并度：$1$。

\bigskip

\noindent\textbf{（3）可分辨粒子}

基态：
\[
\Psi_{\mathrm{gs}}(x_1,x_2)=\psi_1(x_1)\psi_1(x_2),\qquad
E_{\mathrm{gs}}=2E_1,
\]
简并度：$1$。

第一激发态（两种有序组合）：
\[
\Psi_{(1,2)}=\psi_1(x_1)\psi_2(x_2),\qquad
\Psi_{(2,1)}=\psi_2(x_1)\psi_1(x_2),
\]
\[
E_1=E_1+E_2,
\]
简并度：$2$。
\end{solution}
\end{document}